% ============================================================================
% COURS 06 — CORRECTION ET RÉGLAGE DES SYSTÈMES ASSERVIS
% Version propre utilisant les bibliothèques du framework SI
% ============================================================================

\section{Introduction à la correction}\label{sec:corr-introduction}

Un système asservi doit satisfaire simultanément plusieurs exigences : stabilité,
rapidité, précision, limitation du dépassement et robustesse. Le comportement du
procédé seul ne permet pas toujours de respecter le cahier des charges. On ajoute
alors un \textbf{correcteur} dans la chaîne directe afin de modifier la fonction de
transfert en boucle ouverte et, par conséquent, le comportement en boucle fermée.

\TLCorrecteurSerie

Dans la suite, on note :
\[
C(p)\quad\text{le correcteur},\qquad G(p)\quad\text{le procédé},\qquad
H(p)\quad\text{le capteur}.
\]
La fonction de transfert en boucle ouverte est alors
\[
H_{\mathrm{BO}}(p)=C(p)G(p)H(p).
\]
Le correcteur est réglé à partir du cahier des charges et des diagrammes de Bode de
la boucle ouverte.

\subsection{Performances à considérer}

\begin{center}
\renewcommand{\arraystretch}{1.25}
\begin{tabular}{|p{0.22\linewidth}|p{0.68\linewidth}|}
\hline
\textbf{Performance} & \textbf{Caractérisation} \\
\hline
Stabilité & Position des pôles de la FTBF et marges de gain et de phase. \\
\hline
Rapidité & Temps de réponse à 5\,\% ou pulsation de coupure à 0\,dB. \\
\hline
Précision & Erreur statique et classe de la fonction de transfert en boucle ouverte. \\
\hline
Dépassement & Valeur du premier dépassement et amortissement de la réponse. \\
\hline
Robustesse & Conservation des performances malgré les incertitudes du modèle et les perturbations. \\
\hline
\end{tabular}
\end{center}

La stabilité reste prioritaire. L'amélioration d'une performance dégrade souvent
une autre performance : on réalise donc un compromis.

\section{Actions élémentaires de correction}\label{sec:corr-actions}

Les correcteurs usuels sont construits à partir de trois actions élémentaires :
proportionnelle, intégrale et dérivée.

\TLTroisActionsCorrection

\subsection{Action proportionnelle}

\[
u(t)=K\varepsilon(t)
\qquad\Longleftrightarrow\qquad
C(p)=K.
\]
Une augmentation de $K$ augmente le gain de la boucle ouverte. Elle améliore en
général la précision et la rapidité, mais réduit les marges de stabilité.

\subsection{Action intégrale}

\[
u(t)=\frac{1}{T_i}\int_0^t\varepsilon(\tau)\,\mathrm d\tau
\qquad\Longleftrightarrow\qquad
C(p)=\frac{1}{T_i p}.
\]
L'intégrateur augmente la classe de la boucle ouverte. Il permet donc d'annuler
certaines erreurs permanentes. En revanche, il retire $90^\circ$ de phase et peut
fortement dégrader la stabilité.

\subsection{Action dérivée}

\[
u(t)=T_d\frac{\mathrm d\varepsilon(t)}{\mathrm dt}
\qquad\Longleftrightarrow\qquad
C(p)=T_d p.
\]
L'action dérivée anticipe l'évolution de l'écart. Une dérivation idéale est
irréalisable et amplifie fortement le bruit. En pratique, elle est filtrée :
\[
C_D(p)=K\frac{1+T_d p}{1+bT_d p},\qquad 0<b<1.
\]

\section{Correcteur proportionnel P}\label{sec:corr-p}

Le correcteur proportionnel a pour fonction de transfert
\[
\boxed{C(p)=K}.
\]

\TLBodeCorrecteurP[10]

Le diagramme de gain est translaté de $20\log_{10}K$ sans modification de la
phase.

\subsection{Influence sur les performances}

\begin{itemize}
  \item la précision est généralement améliorée ;
  \item la bande passante augmente et le système devient souvent plus rapide ;
  \item les marges de gain et de phase diminuent ;
  \item un gain trop élevé peut entraîner des oscillations ou l'instabilité ;
  \item la saturation de l'actionneur doit être vérifiée.
\end{itemize}

\subsection{Réglage à partir d'une marge de phase}

Pour imposer une marge de phase $M_\varphi$ :
\begin{enumerate}
  \item déterminer la phase souhaitée à la pulsation de coupure :
  \[
  \varphi(\omega_c)=-180^\circ+M_\varphi ;
  \]
  \item relever sur le diagramme non corrigé la pulsation correspondante ;
  \item relever le gain $G_{\mathrm{dB}}(\omega_c)$ ;
  \item choisir
  \[
  20\log_{10}K=-G_{\mathrm{dB}}(\omega_c).
  \]
\end{enumerate}

\section{Correcteurs intégral I et proportionnel intégral PI}\label{sec:corr-pi}

\subsection{Correcteur intégral I}

\[
\boxed{C(p)=\frac{1}{T_i p}}
\]

\TLBodeCorrecteurI[1]

Le gain présente une pente de $-20$\,dB/décade et la phase vaut $-90^\circ$.
Ce correcteur améliore fortement la précision, mais il est rarement utilisé seul
car son effet déstabilisant est important.

\subsection{Correcteur PI}

\[
\boxed{C(p)=K\frac{1+T_i p}{T_i p}}
\]

\TLBodeCorrecteurPI[1][1]

Le correcteur PI se comporte comme un intégrateur aux basses fréquences et comme
un gain proportionnel aux hautes fréquences. Il est particulièrement adapté aux
systèmes qui doivent présenter une erreur statique nulle.

La pulsation de cassure est
\[
\omega_i=\frac{1}{T_i}.
\]
Pour limiter la dégradation de la marge de phase, on place généralement cette
pulsation suffisamment en dessous de la pulsation de coupure souhaitée :
\[
\omega_i\leq \frac{\omega_c}{10}.
\]

\subsection{Compensation d'un pôle dominant}

Si le procédé possède un pôle lent $1+Tp$, on peut choisir $T_i=T$ afin que le
zéro du PI compense ce pôle. Cette compensation doit rester prudente, car le
modèle réel est toujours entaché d'incertitudes.

\section{Correcteur à retard de phase}\label{sec:corr-retard}

Un correcteur à retard de phase s'écrit sous la forme
\[
\boxed{C(p)=K\frac{1+Tp}{1+aTp}},\qquad a>1.
\]

\TLBodeRetardPhase[1][10][1]

Ses pulsations de cassure sont
\[
\omega_p=\frac{1}{aT},\qquad \omega_z=\frac{1}{T},
\qquad \omega_p<\omega_z.
\]
Entre les deux cassures, le gain décroît avec une pente de $-20$\,dB/décade et la
phase devient négative.

Ce correcteur permet d'augmenter le gain aux basses fréquences, donc d'améliorer
la précision, sans modifier fortement le gain aux hautes fréquences. Il tend
toutefois à diminuer la rapidité.

\section{Correcteur à avance de phase}\label{sec:corr-avance}

Un correcteur à avance de phase s'écrit
\[
\boxed{C(p)=K\frac{1+Tp}{1+aTp}},\qquad 0<a<1.
\]

\TLBodeAvancePhase[1][0.1][1]

Les cassures vérifient
\[
\omega_z=\frac{1}{T},\qquad \omega_p=\frac{1}{aT},
\qquad \omega_z<\omega_p.
\]
Le correcteur apporte une phase positive autour de
\[
\omega_m=\frac{1}{T\sqrt a}.
\]
L'avance de phase maximale vaut
\[
\varphi_{\max}=\arcsin\left(\frac{1-a}{1+a}\right).
\]

\subsection{Méthode de réglage}

\begin{enumerate}
  \item déterminer l'avance de phase nécessaire, en conservant une marge de sécurité ;
  \item calculer
  \[
  a=\frac{1-\sin\varphi_{\max}}{1+\sin\varphi_{\max}} ;
  \]
  \item choisir la nouvelle pulsation de coupure au voisinage de $\omega_m$ ;
  \item en déduire
  \[
  T=\frac{1}{\omega_m\sqrt a} ;
  \]
  \item régler $K$ pour que le gain corrigé soit nul à la pulsation de coupure.
\end{enumerate}

Le correcteur à avance de phase améliore les marges de stabilité et augmente la
bande passante. Il améliore donc généralement la rapidité.

\section{Correcteur PID}\label{sec:corr-pid}

Une forme classique du correcteur PID est
\[
\boxed{C(p)=K\left(1+\frac{1}{T_i p}+T_d p\right)}.
\]
Une forme physiquement réalisable utilise une dérivée filtrée :
\[
\boxed{C(p)=K\frac{1+T_i p}{T_i p}\frac{1+T_d p}{1+bT_d p}},
\qquad 0<b<1.
\]

\TLBodeCorrecteurPID[1][10][1][0.1]

\begin{itemize}
  \item l'action proportionnelle règle le niveau global de commande ;
  \item l'action intégrale améliore la précision ;
  \item l'action dérivée filtrée améliore la stabilité et la rapidité.
\end{itemize}

Le PID offre de nombreuses possibilités de réglage, mais il est plus sensible aux
incertitudes, au bruit et aux saturations. En pratique, on utilise souvent un PI
ou un PID filtré avec un dispositif anti-saturation de l'intégrateur.

\section{Bilan et architectures de correction}\label{sec:corr-bilan}

\begin{center}
\renewcommand{\arraystretch}{1.25}
\begin{tabular}{|p{0.14\linewidth}|c|c|c|p{0.28\linewidth}|}
\hline
\textbf{Correcteur} & \textbf{Précision} & \textbf{Rapidité} & \textbf{Stabilité} & \textbf{Remarque} \\
\hline
P & $+$ & $+$ & $-$ & Réglage simple, attention à la saturation. \\
\hline
I & $++$ & $-$ & $--$ & Rarement employé seul. \\
\hline
PI & $++$ & $0/-$ & $-$ & Très courant pour annuler l'erreur statique. \\
\hline
Retard & $+$ & $-$ & $0/-$ & Améliore le gain en basse fréquence. \\
\hline
Avance & $0$ & $+$ & $+$ & Augmente la marge de phase et la bande passante. \\
\hline
PID & $++$ & $+$ & $+$ & Performant mais plus délicat à régler. \\
\hline
\end{tabular}
\end{center}

\subsection{Boucles imbriquées}

Pour un axe motorisé, on utilise fréquemment une boucle interne de vitesse et une
boucle externe de position. La boucle interne doit être nettement plus rapide que
la boucle externe.

\TLBouclesImbriquees

Une règle pratique consiste à séparer les bandes passantes d'un facteur compris
entre 5 et 10. On règle d'abord la boucle interne, puis la boucle externe en
considérant la boucle interne fermée comme un sous-système.

\subsection{Démarche générale de synthèse}

\begin{enumerate}
  \item établir et valider le modèle du procédé ;
  \item analyser la boucle ouverte non corrigée ;
  \item traduire le cahier des charges en contraintes fréquentielles ;
  \item choisir la structure du correcteur ;
  \item calculer ses paramètres ;
  \item vérifier stabilité, précision, rapidité et saturation ;
  \item valider par simulation puis expérimentalement ;
  \item ajuster le réglage en conservant des marges de robustesse.
\end{enumerate}

\section{À retenir}\label{sec:corr-retenir}

\begin{itemize}
  \item Le réglage s'effectue principalement sur la fonction de transfert en boucle ouverte.
  \item Le gain proportionnel améliore précision et rapidité, mais réduit la stabilité.
  \item L'intégrateur augmente la classe et annule certaines erreurs permanentes.
  \item L'avance de phase améliore la marge de phase et la rapidité.
  \item Le retard de phase améliore surtout la précision aux basses fréquences.
  \item Un PID combine les trois actions, avec une dérivée nécessairement filtrée.
  \item La stabilité, les saturations et la robustesse doivent toujours être vérifiées sur le système réel.
\end{itemize}
